%%%%%%%%%%%% ------------ chapter header ------------  %%%%%%%%%%%%
\chapterhead
% chapter number, title
{序論}
% chapter abst
{本章では，本研究の背景と目的および，本論文の内容構成について述べる．}

%%%%%%%%%%%% ------------ next page ------------  %%%%%%%%%%%%

% \section{}で小見出しを作る
% 　章ごとにファイルを分けている

\section{研究背景}
研究分野全体の現状や社会的背景について述べる．
なぜこの研究テーマが重要なのか，どのような課題が存在するのかを説明する．

% 文章作成時の注意点
% \parで段落を変える
% ","と"，"　"."と"．"の違いに気をつける
% 表記，特に英数字を気をつける(例: M5Stack，m5stackなど)
% 文末は揃える
% \cite{タグ名}で引用　参考文献の欄に詳細を書く

% 箇条書きの例
\begin{itemize}
  \item 現在の研究分野における課題
  \item 社会的ニーズや技術的背景
  \item 問題点の具体例
\end{itemize}

\newpage

\section{研究目的}
本研究の目的を明確に述べる．
背景で述べた課題に対して，本研究がどのようなアプローチで取り組むのかを説明する．

本研究では，〜〜システムを提案する．
このシステムにより，〜〜を実現することを目的とする．

\newpage

\section{本論文の構成}
本論文の構成は以下の通りである．
各章で何を述べるのかを簡潔に説明する．

第2章では，関連研究について述べる．
第3章では，提案手法について述べる．
第4章では，実験方法と実験条件について述べる．
第5章では，実験結果を示し，考察を行う．
第6章では，本論文の結論を述べる．
\newpage